\section[Sampling with fair bits]{Bounded categorical sampling from fair bits}
\label{sec:op3-fair-bit-sampling}

\paragraph{Purpose and model.}
The reference confidence audit uses a common-denominator integer sampler,
with its integer operations explicitly identified as reference work.
For the supplied constructor we now give a sampler using only exact-real
word arithmetic, comparisons, integer indices and independent fair bits.
No uniformly random real, common-denominator integer sampler, or minimum
positive probability is imported. Dyadic arithmetic uses the same
exact-real word model as the supplied solver. The result below charges
those operations and makes no assertion about their bit complexity.

\begin{proposition}[Capped fair-bit categorical sampling; \statusproved]
\label{prop:op3-fair-bit-categorical}
Given positive weights $a_1,\ldots,a_r$, store their prefix sums in
$O(r)$ work and words. For any integer $B\geq0$, a draw uses at most
$B$ fair bits and $O((B+1)\log(2+r))$ work. It returns a category or
\textnormal{FAIL}. It couples to an exact draw with probabilities
$a_i/\sum_j a_j$, agrees with that draw whenever it returns a category,
and satisfies
\[
 \Pr\{\textnormal{FAIL}\}\leq (r-1)2^{-B}.
\]
With no cap, expected work is $O(\log^2(2+r))$.
All prefix construction, searches, arithmetic, fair bits, temporary
state and output records are charged.

For at most $N$ draws whose distributions may depend on previous
successful outputs, suppose every distribution has at most $R$
categories. If $N(R-1)2^{-B}\leq p$, the joint coupling fails with
probability at most $p$. The bounded sampling work is
$O(N(B+1)\log(2+R))$, in addition to all distribution-assembly work.
For $R=1$ no random bit is necessary.
\end{proposition}
\begin{proof}
Write $A=\sum_i a_i$ and $s_i=\sum_{j\leq i}a_j$. The first $k$
fair bits define an integer $z_k\in\{0,\ldots,2^k-1\}$ and the
dyadic interval
\[
 I_k=[z_k/2^k,(z_k+1)/2^k).
\]
Start with $I_0=[0,1)$. Find the first $i$ with
$2^k s_i>z_k A$ by binary search. If
$(z_k+1)A\leq2^k s_i$, the whole interval is in category $i$;
return it. Otherwise read another fair bit unless $k=B$, in which
case return \textnormal{FAIL}. The strict comparison at the lower
endpoint and nonstrict comparison at the upper endpoint handle dyadic
CDF boundaries without introducing an error.

Imagine completing each draw's bit sequence indefinitely to obtain a
uniform point $U\in[0,1)$. Every interval inspected contains that point.
When the procedure returns, its category is exactly the category of $U$.
At depth $k$, an unresolved dyadic interval must contain a category
boundary strictly in its interior. There are $r-1$ boundaries and
therefore at most $r-1$ unresolved intervals among the $2^k$ equally
likely intervals. This proves the failure bound. Binary-expansion
ambiguities at individual points have probability zero.

If $T$ is the number of bits required without a cap, then
$\Pr\{T>k\}\leq\min\{1,(r-1)2^{-k}\}$. Summing this tail gives
$\mathbb E T=O(\log(2+r))$, while each interval test costs
$O(\log(2+r))$. The computation retains only its current integer
numerator, dyadic scale and binary-search indices; it never stores all
old prefixes. Counting even all created scalar temporaries is bounded
by the stated word work.

For adaptive draws, couple the two processes until the first failure.
Conditional on any shared previous history, the next exact draw uses
fresh fair bits, and the conditional failure probability has the same
uniform bound $(R-1)2^{-B}$. A union bound proves the joint claim.
A sufficient $B$ can be found by doubling an integer scale until it
exceeds $N(R-1)/p$, charging $O(1+\log(NR/p))$ additional operations.
The cases $N=0$ or $R=1$ are handled directly.
\end{proof}

\paragraph{Use the coupling, not conditioning on successful draws.}
For weights $(1,2)$ and cap $B=1$, the lower dyadic half is unresolved,
while the upper half returns the second category. Thus conditional on
nonfailure, that category has probability one, rather than $2/3$.
When applying a concentration theorem, analyze the exact coupled
sampling process and add the capped failure probability. Do not assert
that the accepted conditional distribution is unchanged.

\paragraph{Verification.}
The implementation \texttt{fair\_bit\_categorical.py} independently
checks terminal dyadic intervals against exact category intersections,
exhaustive short-bit distribution and failure bounds, adaptive pairs,
and explicitly chosen rare-event strings. Its source hash and full
counts are saved in \texttt{FAIR\_BIT\_CATEGORICAL\_AUDIT.json}.
The full audit checks 140,860 successful terminal intervals, 879
ambiguous terminal intervals, 209 complete distribution/failure grids,
and seven adaptive two-draw union bounds. Explicit rare-event strings
reach probability $2^{-1024}$ without rounding. Seeded sample frequencies
are not used to prove the probability guarantee.
This supplies a sampling primitive; the complete resistance estimator,
core sparsifier and local OP3 algorithm are not asserted by it.
